\subsection{Bilinear Maps and the Tensor Product Universal Property}

% Role: bilinear maps, the tensor-product universal property, its quotient
% construction, uniqueness, and the bilinear-linear correspondence.

Tensor products are the device that turns bilinear data into linear data. The guiding rule is the same one already used for quotient spaces in Theorem \ref{thm:quotient-universal-property}: a universal object is recognised by the unique way every compatible map factors through it.

\begin{definition} \label{def:bilinear-map} \hlt{Bilinear Maps}\\
Let $\mathcal{U},\mathcal{V},\mathcal{W}$ be $\mathbb{F}$-vector spaces and $f:\mathcal{U}\times\mathcal{V}\to\mathcal{W}$.
\begin{enumerate}[label=(\roman*)]
\setlength{\itemsep}{0pt}
\item $f$ is \hlt{bilinear} if $\forall r,s\in\mathbb{F}$, $\mathbf{u},\mathbf{u}'\in\mathcal{U}$, $\mathbf{v},\mathbf{v}'\in\mathcal{V}$, $f(r\mathbf{u}+s\mathbf{u}',\mathbf{v})=r f(\mathbf{u},\mathbf{v})+s f(\mathbf{u}',\mathbf{v})$\\
and $f(\mathbf{u},r\mathbf{v}+s\mathbf{v}')=r f(\mathbf{u},\mathbf{v})+s f(\mathbf{u},\mathbf{v}')$.
\item $\operatorname{Bil}_{\mathbb{F}}(\mathcal{U},\mathcal{V};\mathcal{W})\defeq\{f:\mathcal{U}\times\mathcal{V}\to\mathcal{W}:f\text{ bilinear}\}$.
\item If $\mathcal{W}=\mathbb{F}$, then $\operatorname{Bil}_{\mathbb{F}}(\mathcal{U},\mathcal{V};\mathbb{F})$ is the space of \hlt{bilinear forms} from Definition \ref{def:bilinear-form-and-matrix}.
\end{enumerate}
\end{definition}

\begin{example} \label{ex:basic-bilinear-maps} \hlt{Basic Bilinear Maps}
\begin{enumerate}[label=(\roman*)]
\setlength{\itemsep}{0pt}
\item If $A$ is an $\mathbb{F}$-algebra, then $\mu:A\times A\to A$, $\mu(a,b)\defeq ab$, is bilinear.
\item Evaluation $\operatorname{ev}:\LL(\mathcal{V},\mathcal{W})\times\mathcal{V}\to\mathcal{W}$, $\operatorname{ev}(T,\mathbf{v})\defeq T\mathbf{v}$, is bilinear.
\item If $\varphi\in\mathcal{V}'$ and $\psi\in\mathcal{W}'$, then $B_{\varphi,\psi}:\mathcal{V}\times\mathcal{W}\to\mathbb{F}$, $B_{\varphi,\psi}(\mathbf{v},\mathbf{w})\defeq\varphi(\mathbf{v})\psi(\mathbf{w})$, is bilinear.
\end{enumerate}
\end{example}

\begin{definition} \label{def:tensor-product-universal-property} \hlt{Tensor Product by Universal Property}\\
Let $\mathcal{U},\mathcal{V}$ be $\mathbb{F}$-vector spaces.
\begin{enumerate}[label=(\roman*)]
\setlength{\itemsep}{0pt}
\item A pair $(\mathcal{T},t)$ is a \hlt{tensor product} of $\mathcal{U}$ and $\mathcal{V}$ if $\mathcal{T}$ is an $\mathbb{F}$-vector space, $t\in\operatorname{Bil}_{\mathbb{F}}(\mathcal{U},\mathcal{V};\mathcal{T})$, and for every $\mathbb{F}$-vector space $\mathcal{W}$ and $f\in\operatorname{Bil}_{\mathbb{F}}(\mathcal{U},\mathcal{V};\mathcal{W})$, $\exists!\widetilde{f}\in\LL(\mathcal{T},\mathcal{W})$ such that $f=\widetilde{f}\circ t$.
\item The map $\widetilde{f}$ is the \hlt{mediating map} for $f$.
\item Write $\mathcal{T}=\mathcal{U}\otimes\mathcal{V}$ and $\mathbf{u}\otimes\mathbf{v}\defeq t(\mathbf{u},\mathbf{v})$; such tensors are \hlt{decomposable}.
\end{enumerate}
\end{definition}

\begin{figure}[H]
  \centering
  % Source note: Roman, Advanced Linear Algebra, Ch. 14, Fig. 14.4.
% Redrawn as a schematic universal-property diagram for these notes.
\begin{tikzpicture}[
  >=Stealth,
  every node/.style={font=\small},
  arrow/.style={->, line width=0.55pt},
  factor/.style={->, dashed, line width=0.55pt},
  note/.style={font=\scriptsize, text=gray!70!black}
]
  \node (uv) at (0,1.25) {$\mathcal{U}\times\mathcal{V}$};
  \node (t) at (3.6,1.25) {$\mathcal{U}\otimes\mathcal{V}$};
  \node (w) at (3.6,-0.8) {$\mathcal{W}$};

  \draw[arrow] (uv) -- node[above, note] {$t$ bilinear} (t);
  \draw[arrow] (uv) -- node[below left, note] {$f$ bilinear} (w);
  \draw[factor] (t) -- node[right, note] {$\widetilde{f}$ linear} (w);
\end{tikzpicture}

  \caption{A tensor product is universal for bilinear maps out of $\mathcal{U}\times\mathcal{V}$.}
  \label{fig:tensor-universal-property}
\end{figure}

\begin{theorem} \label{thm:tensor-product-quotient-construction} \hlt{Coordinate-Free Quotient Construction}\\
Let $\mathcal{U},\mathcal{V}$ be $\mathbb{F}$-vector spaces and let $\mathbb{F}^{(\mathcal{U}\times\mathcal{V})}$ have basis $\{[\mathbf{u},\mathbf{v}]:(\mathbf{u},\mathbf{v})\in\mathcal{U}\times\mathcal{V}\}$.\\
Let $\mathcal{R}$ be the span of all $r[\mathbf{u}_1,\mathbf{v}]+s[\mathbf{u}_2,\mathbf{v}]-[r\mathbf{u}_1+s\mathbf{u}_2,\mathbf{v}]$; and $r[\mathbf{u},\mathbf{v}_1]+s[\mathbf{u},\mathbf{v}_2]-[\mathbf{u},r\mathbf{v}_1+s\mathbf{v}_2]$, where $r,s\in\mathbb{F}$, $\mathbf{u},\mathbf{u}_1,\mathbf{u}_2\in\mathcal{U}$, and $\mathbf{v},\mathbf{v}_1,\mathbf{v}_2\in\mathcal{V}$. Then $\mathcal{U}\otimes\mathcal{V}\defeq \mathbb{F}^{(\mathcal{U}\times\mathcal{V})}/\mathcal{R}$, $t(\mathbf{u},\mathbf{v})=[\mathbf{u},\mathbf{v}]+\mathcal{R}$ is a tensor product of $\mathcal{U}$ and $\mathcal{V}$.
\begin{proof} Let $\pi:\mathbb{F}^{(\mathcal{U}\times\mathcal{V})}\to\mathcal{U}\otimes\mathcal{V}$ be the quotient map.\\
Since the generators of $\mathcal{R}$ vanish in the quotient, $t(\mathbf{u},\mathbf{v})=\pi[\mathbf{u},\mathbf{v}]$ is bilinear.\\
Let $f\in\operatorname{Bil}_{\mathbb{F}}(\mathcal{U},\mathcal{V};\mathcal{W})$ and define $\alpha:\mathbb{F}^{(\mathcal{U}\times\mathcal{V})}\to\mathcal{W}$ by $\alpha[\mathbf{u},\mathbf{v}]\defeq f(\mathbf{u},\mathbf{v})$.\\
Bilinearity gives $\mathcal{R}\subseteq\ker\alpha$, so Theorem \ref{thm:quotient-universal-property} gives a unique $\widetilde{f}\in\LL(\mathcal{U}\otimes\mathcal{V},\mathcal{W})$ with $\widetilde{f}\circ \pi=\alpha$. Hence $\widetilde{f}\circ t=f$.\\
If $L\circ t=f$, then $L$ and $\widetilde{f}$ agree on the spanning set $\{[\mathbf{u},\mathbf{v}]+\mathcal{R}\}$, so $L=\widetilde{f}$.\\
\hlt{Structural Commentary}: The quotient kills precisely the failure of the raw ordered-pair basis to be bilinear.\\
\hlt{Key Insight}: A tensor product is a free bilinear map, followed by no extra identifications.
\end{proof}
\end{theorem}

\begin{figure}[H]
  \centering
  % Source note: Roman, Advanced Linear Algebra, Ch. 14, Fig. 14.5.
% Redrawn to show the free-vector-space quotient construction.
\begin{tikzpicture}[
  >=Stealth,
  x=0.9cm,
  y=0.9cm,
  every node/.style={font=\small},
  arrow/.style={->, line width=0.65pt, black},
  aux/.style={->, dashed, line width=0.65pt, black},
  factor/.style={->, dashed, line width=0.7pt, black},
  note/.style={font=\scriptsize, text=black, fill=white, inner sep=0.8pt}
]
  \node (uv) at (0,0.45) {$\mathcal{U}\times\mathcal{V}$};
  \node (free) at (3.2,0.45) {$\mathbb{F}^{(\mathcal{U}\times\mathcal{V})}$};
  \node (quot) at (6.55,0.45) {$\mathcal{U}\otimes\mathcal{V}$};
  \node (w) at (3.2,-0.58) {$\mathcal{W}$};

  \draw[arrow] (uv) -- node[above, note] {$j$} (free);
  \draw[arrow] (free) -- node[above, note, pos=0.55] {$\pi$} (quot);
  \draw[arrow] (uv.north east) to[out=30,in=150]
    node[above, note] {$t=\pi j$} (quot.north west);
  \draw[arrow] (uv) -- node[below left, note] {$f$} (w);
  \draw[aux] (free) -- node[right, note, xshift=2pt] {$\alpha$} (w);
  \draw[factor] (quot) -- node[right, note, xshift=8pt] {$\widetilde{f}$} (w);
\end{tikzpicture}

  \caption{The free linear map $\alpha$ kills $\mathcal{R}$, so it descends
  uniquely to $\widetilde{f}:\mathcal{U}\otimes\mathcal{V}\to\mathcal{W}$.}
  \label{fig:tensor-quotient-construction-universal}
\end{figure}

\begin{corollary} \label{cor:tensor-product-unique-up-to-isomorphism} \hlt{Tensor Products Are Unique up to Isomorphism}\\
If $(\mathcal{T},t)$ and $(\mathcal{S},s)$ are tensor products of $\mathcal{U}$ and $\mathcal{V}$,\\
then there is a unique isomorphism $\Phi:\mathcal{T}\to\mathcal{S}$ satisfying $\Phi(t(\mathbf{u},\mathbf{v}))=s(\mathbf{u},\mathbf{v})$, where $\mathbf{u}\in\mathcal{U},\ \mathbf{v}\in\mathcal{V}$.
\begin{proof}
Apply universal property of $(\mathcal{T},t)$ to bilinear map $s$ and universal property of $(\mathcal{S},s)$ to bilinear map $t$.\\
This gives linear maps $\Phi:\mathcal{T}\to\mathcal{S}$ and $\Psi:\mathcal{S}\to\mathcal{T}$. The composites $\Psi\Phi$ and $I_{\mathcal{T}}$ are both mediating maps for $t$, while $\Phi\Psi$ and $I_{\mathcal{S}}$ are both mediating maps for $s$. Uniqueness forces $\Psi\Phi=I_{\mathcal{T}}$ and $\Phi\Psi=I_{\mathcal{S}}$.\\
\hlt{Structural Commentary}: Universal objects are unique up to canonical isomorphism.\\
\hlt{Key Insight}: Two universal factorisations force the comparison maps to be inverse.
\end{proof}
\end{corollary}

\begin{theorem} \label{thm:tensor-linearizes-bilinear-maps} \hlt{Bilinearity Equals Linearity on a Tensor Product}\\
For $\mathbb{F}$-vector spaces $\mathcal{U},\mathcal{V},\mathcal{W}$, the map $\Phi:\operatorname{Bil}_{\mathbb{F}}(\mathcal{U},\mathcal{V};\mathcal{W})\to\LL(\mathcal{U}\otimes\mathcal{V},\mathcal{W})$, $\Phi(f)=\widetilde{f}$,\\
where $\widetilde{f}(\mathbf{u}\otimes\mathbf{v})=f(\mathbf{u},\mathbf{v})$, is an isomorphism.
\begin{proof}
The universal property gives existence and uniqueness of $\widetilde{f}$.\\
Define $\Psi:\LL(\mathcal{U}\otimes\mathcal{V},\mathcal{W})\to\operatorname{Bil}_{\mathbb{F}}(\mathcal{U},\mathcal{V};\mathcal{W})$ by $\Psi(L)\defeq L\circ t$.\\
Then $\Psi(\Phi(f))=\widetilde{f}\circ t=f$, and $\Phi(\Psi(L))=L$ by uniqueness of the mediating map for $L\circ t$. Hence $\Phi$ is bijective.\\
For $a\in\mathbb{F}$ and $f,g$ bilinear, $\Phi(af+g)(\mathbf{u}\otimes\mathbf{v})=af(\mathbf{u},\mathbf{v})+g(\mathbf{u},\mathbf{v})=(a\Phi(f)+\Phi(g))(\mathbf{u}\otimes\mathbf{v})$, so $\Phi$ is linear.\\
\hlt{Structural Commentary}: Tensor products package bilinear maps as ordinary linear maps.\\
\hlt{Key Insight}: To define a linear map out of $\mathcal{U}\otimes\mathcal{V}$, define a bilinear rule on $\mathcal{U}\times\mathcal{V}$.
\end{proof}
\end{theorem}

\begin{method} \label{meth:defining-linear-maps-on-tensor-products} \hlt{Defining Maps on Tensor Products}\\
To define $L\in\LL(\mathcal{U}\otimes\mathcal{V},\mathcal{W})$:
\begin{enumerate}[label=(\roman*)]
\setlength{\itemsep}{0pt}
\item define a candidate rule $f:\mathcal{U}\times\mathcal{V}\to\mathcal{W}$ by $f(\mathbf{u},\mathbf{v})$;
\item prove $f\in\operatorname{Bil}_{\mathbb{F}}(\mathcal{U},\mathcal{V};\mathcal{W})$;
\item set $L\defeq\Phi(f)$ from Theorem \ref{thm:tensor-linearizes-bilinear-maps}, so $L(\mathbf{u}\otimes\mathbf{v})=f(\mathbf{u},\mathbf{v})$.
\end{enumerate}
Do not assign arbitrary values on decomposable tensors: they span $\mathcal{U}\otimes\mathcal{V}$ but are usually linearly dependent.
\end{method}
